\documentclass[border=5pt]{standalone}
\usepackage[T1]{fontenc}
\usepackage{tikz}
\usepackage{luacode}

\definecolor{garnet}{HTML}{73000A}
\definecolor{atlantic}{HTML}{466A9F}
\definecolor{congaree}{HTML}{1F414D}
\definecolor{warmgrey}{HTML}{676156}
\definecolor{bk70}{HTML}{5C5C5C}
\definecolor{bk50}{HTML}{A2A2A2}
\definecolor{bk30}{HTML}{C7C7C7}

\begin{document}
\begin{tikzpicture}

\begin{luacode}
-- ================================================================
-- Two-lens Snell's-law ray tracer
-- ================================================================

-- ── Adjustable parameters ──
local lens1_x  = 3.0     -- center x of lens 1
local lens2_x  = 5.5     -- center x of lens 2
local lens1_R  = 4.0     -- radius of curvature (smaller = more curved)
local lens2_R  = 5.0
local lens1_h  = 1.3     -- half-aperture
local lens2_h  = 1.3
local n_glass  = 1.52
local n_air    = 1.0
local sensor_x = 8.0
local sensor_h = 1.5

local xmin, xmax = -0.5, 8.8
local ymin, ymax = -2.0, 2.0

-- ── Helpers ──
local function sag(R, h) return R - math.sqrt(math.max(R*R - h*h, 0)) end

-- Circle-ray intersection.  Returns (t, hx, hy, nx, ny).
-- pick = "near" or "far" selects which root.
local function isect(ox, oy, dx, dy, cx, cy, r, pick)
    local ex, ey = ox - cx, oy - cy
    local a = dx*dx + dy*dy
    local b = 2*(ex*dx + ey*dy)
    local c = ex*ex + ey*ey - r*r
    local disc = b*b - 4*a*c
    if disc < 0 then return nil end
    local sq = math.sqrt(disc)
    local t1 = (-b - sq)/(2*a)
    local t2 = (-b + sq)/(2*a)
    local t
    if pick == "far" then
        t = (t2 > 1e-4) and t2 or ((t1 > 1e-4) and t1 or nil)
    else
        t = (t1 > 1e-4) and t1 or ((t2 > 1e-4) and t2 or nil)
    end
    if not t then return nil end
    local hx, hy = ox + t*dx, oy + t*dy
    return t, hx, hy, (hx-cx)/r, (hy-cy)/r
end

-- Vector-form Snell refraction.  Returns (dx', dy') or nil (TIR).
local function refract(dx, dy, nx, ny, n1, n2)
    local d = dx*nx + dy*ny
    if d > 0 then nx, ny, d = -nx, -ny, -d end   -- flip normal to face ray
    local r = n1/n2
    local cos_i = -d
    local sin2t = r*r*(1 - cos_i*cos_i)
    if sin2t > 1 then return nil end
    local cos_t = math.sqrt(1 - sin2t)
    local rx = r*dx + (r*cos_i - cos_t)*nx
    local ry = r*dy + (r*cos_i - cos_t)*ny
    local l = math.sqrt(rx*rx + ry*ry)
    return rx/l, ry/l
end

-- ── Biconvex lens geometry ──
-- Left surface: part of circle centered at  cx_L = lx + sqrt(R²-h²)
--   surface is the LEFT side of that circle (ray enters here)
-- Right surface: part of circle centered at cx_R = lx - sqrt(R²-h²)
--   surface is the RIGHT side of that circle (ray exits here)

local function lens_circles(lx, R, h)
    local d = math.sqrt(math.max(R*R - h*h, 0))
    return lx + d, lx - d   -- cx_left_surf, cx_right_surf
end

-- ── Draw a biconvex lens ──
local function draw_lens(lx, R, h)
    local s = sag(R, h)
    local cx_L, cx_R = lens_circles(lx, R, h)
    local N = 60

    -- Left surface arc: x = cx_L - sqrt(R²-y²)  for y in [-h, h]
    -- Right surface arc: x = cx_R + sqrt(R²-y²)  for y in [-h, h]
    local left, right = {}, {}
    for i = 0, N do
        local y = -h + 2*h*(i/N)
        local q = math.sqrt(math.max(R*R - y*y, 0))
        table.insert(left,  {cx_L - q, y})
        table.insert(right, {cx_R + q, y})
    end

    -- Fill
    local p = "\\fill[atlantic!10, opacity=0.85] "
    for i, v in ipairs(left) do
        p = p .. (i==1 and "" or " -- ") .. string.format("(%.4f,%.4f)", v[1], v[2])
    end
    for i = #right, 1, -1 do
        p = p .. string.format(" -- (%.4f,%.4f)", right[i][1], right[i][2])
    end
    tex.sprint(p .. " -- cycle;")

    -- Outlines
    for _, arc in ipairs({left, right}) do
        local s = "\\draw[atlantic!60, thin] "
        for i, v in ipairs(arc) do
            s = s .. (i==1 and "" or " -- ") .. string.format("(%.4f,%.4f)", v[1], v[2])
        end
        tex.sprint(s .. ";")
    end
    -- Top/bottom edges
    tex.sprint(string.format(
        "\\draw[atlantic!60, thin] (%.4f,%.4f)--(%.4f,%.4f);",
        left[1][1], left[1][2], right[1][1], right[1][2]))
    tex.sprint(string.format(
        "\\draw[atlantic!60, thin] (%.4f,%.4f)--(%.4f,%.4f);",
        left[#left][1], left[#left][2], right[#right][1], right[#right][2]))
end

-- ── Trace one ray through both lenses ──
local function trace(ox, oy, dx, dy)
    local pts = {{ox, oy}}
    local cx, cy, cdx, cdy = ox, oy, dx, dy

    local lenses = {
        {x=lens1_x, R=lens1_R, h=lens1_h},
        {x=lens2_x, R=lens2_R, h=lens2_h},
    }

    for _, L in ipairs(lenses) do
        local cxL, cxR = lens_circles(L.x, L.R, L.h)

        -- 1) Hit left surface (near side of circle centered at cxL)
        local t, hx, hy, nx, ny = isect(cx, cy, cdx, cdy, cxL, 0, L.R, "near")
        if t and math.abs(hy) <= L.h then
            table.insert(pts, {hx, hy})
            local ndx, ndy = refract(cdx, cdy, nx, ny, n_air, n_glass)
            if not ndx then break end
            cdx, cdy = ndx, ndy
            cx, cy = hx, hy

            -- 2) Hit right surface (far side of circle centered at cxR)
            t, hx, hy, nx, ny = isect(cx, cy, cdx, cdy, cxR, 0, L.R, "far")
            if t and math.abs(hy) <= L.h then
                table.insert(pts, {hx, hy})
                ndx, ndy = refract(cdx, cdy, nx, ny, n_glass, n_air)
                if not ndx then break end
                cdx, cdy = ndx, ndy
                cx, cy = hx, hy
            end
        end
        -- if missed, ray continues undeviated
    end

    -- Terminate at sensor plane
    if math.abs(cdx) > 1e-6 then
        local ts = (sensor_x - cx) / cdx
        if ts > 0 then
            table.insert(pts, {sensor_x, cy + ts*cdy})
        end
    end
    return pts
end

-- ══════════════════════════════════
-- DRAW
-- ══════════════════════════════════

-- Optical axis
tex.sprint(string.format(
    "\\draw[warmgrey, dashed, thin] (%.2f,0)--(%.2f,0);", xmin, xmax))

-- Lenses
draw_lens(lens1_x, lens1_R, lens1_h)
draw_lens(lens2_x, lens2_R, lens2_h)

-- CMOS sensor
tex.sprint(string.format(
    "\\fill[congaree!80] (%.2f,%.2f) rectangle (%.2f,%.2f);",
    sensor_x, -sensor_h, sensor_x+0.12, sensor_h))
for i = 0, 15 do
    local y = -sensor_h + i*(2*sensor_h/16)
    tex.sprint(string.format(
        "\\draw[congaree!40, very thin] (%.2f,%.4f)--(%.2f,%.4f);",
        sensor_x, y, sensor_x+0.12, y))
end
tex.sprint(string.format(
    "\\draw[congaree, thick] (%.2f,%.2f) rectangle (%.2f,%.2f);",
    sensor_x, -sensor_h, sensor_x+0.12, sensor_h))

-- ── Parallel rays (garnet) ──
for i = 0, 12 do
    local y0 = -1.2 + 2.4*(i/12)
    local pts = trace(xmin, y0, 1, 0)
    for j = 1, #pts-1 do
        tex.sprint(string.format(
            "\\draw[garnet, thin, opacity=0.6] (%.4f,%.4f)--(%.4f,%.4f);",
            pts[j][1], pts[j][2], pts[j+1][1], pts[j+1][2]))
    end
end

-- ── Point-source fan (atlantic) ──
local sx, sy = -0.3, 0.9
for i = 0, 10 do
    local ty = -lens1_h*0.85 + 2*lens1_h*0.85*(i/10)
    local ddx, ddy = lens1_x - sx, ty - sy
    local l = math.sqrt(ddx*ddx + ddy*ddy)
    local pts = trace(sx, sy, ddx/l, ddy/l)
    for j = 1, #pts-1 do
        tex.sprint(string.format(
            "\\draw[atlantic, thin, opacity=0.45] (%.4f,%.4f)--(%.4f,%.4f);",
            pts[j][1], pts[j][2], pts[j+1][1], pts[j+1][2]))
    end
end
tex.sprint(string.format("\\fill[atlantic] (%.2f,%.2f) circle(0.06);", sx, sy))

-- ── Labels ──
tex.sprint(string.format(
    "\\node[font=\\sffamily\\footnotesize,text=atlantic,anchor=south] at (%.2f,%.2f) {Lens 1};",
    lens1_x, lens1_h+0.12))
tex.sprint(string.format(
    "\\node[font=\\sffamily\\footnotesize,text=atlantic,anchor=south] at (%.2f,%.2f) {Lens 2};",
    lens2_x, lens2_h+0.12))
tex.sprint(string.format(
    "\\node[font=\\sffamily\\footnotesize,text=congaree,anchor=south] at (%.2f,%.2f) {CMOS};",
    sensor_x+0.06, sensor_h+0.12))

-- ── Parameter readout ──
tex.sprint(string.format(
    "\\node[font=\\sffamily\\scriptsize,text=warmgrey,anchor=north west] at (%.2f,%.2f) {$R_1=%.1f$ \\quad $R_2=%.1f$ \\quad $d=%.1f$};",
    xmin+0.1, ymin+0.05, lens1_R, lens2_R, lens2_x-lens1_x))

\end{luacode}

\useasboundingbox (-0.5,-2.0) rectangle (8.8,2.0);

\end{tikzpicture}
\end{document}
